\documentclass[10pt, aspectratio=169]{beamer}
% Preámbulo modular para presentaciones Beamer (Modelación Estadística)
% Diseño y paleta institucional del Tecnológico de Monterrey

\documentclass[aspectratio=169,xcolor={dvipsnames,table}]{beamer}

\usepackage[utf8]{inputenc}
\usepackage[spanish,mexico]{babel}
\usepackage{amsmath,amssymb,amsfonts,mathtools}
\usepackage{graphicx}
\usepackage{listings}
\usepackage{booktabs}
\usepackage{tikz}
\usetikzlibrary{matrix,arrows,decorations.pathmorphing,shapes.geometric,positioning}

% Paleta Institucional Tecnológico de Monterrey
\definecolor{TecAzul}{HTML}{0039A6}       % Azul TEC: color primario institucional
\definecolor{TecAzulOscuro}{HTML}{1A2E51} % Azul marino oscuro
\definecolor{TecRojo}{HTML}{EC2661}       % Rojo de acento (paleta miTec)
\definecolor{TecGris}{HTML}{646464}       % Gris neutro para comentarios y subtítulos
\definecolor{TecFondoBloque}{HTML}{F4F6F9}% Fondo suave para bloques

% Configuración de Tema Beamer
\usetheme{Boadilla}
\usecolortheme[named=TecAzulOscuro]{structure}
\setbeamercolor{alerted text}{fg=TecRojo}
\setbeamercolor{palette primary}{bg=TecAzulOscuro,fg=white}
\setbeamercolor{palette secondary}{bg=TecAzul,fg=white}
\setbeamercolor{palette tertiary}{bg=TecAzulOscuro!80!black,fg=white}
\setbeamercolor{title}{fg=TecAzulOscuro,bg=TecFondoBloque}
\setbeamercolor{frametitle}{fg=TecAzulOscuro,bg=TecFondoBloque}

% Bloques personalizados elegantes
\setbeamercolor{block title}{bg=TecAzulOscuro,fg=white}
\setbeamercolor{block body}{bg=TecFondoBloque,fg=black}
\setbeamercolor{block title alerted}{bg=TecRojo,fg=white}
\setbeamercolor{block body alerted}{bg=TecRojo!8,fg=black}
\setbeamercolor{block title example}{bg=TecAzul,fg=white}
\setbeamercolor{block body example}{bg=TecAzul!8,fg=black}

% Redondear bordes y añadir sombra sutil
\setbeamertemplate{blocks}[rounded][shadow=true]
\setbeamertemplate{navigation symbols}{}

% Pie de página limpio con número de diapositiva
\setbeamertemplate{footline}{
  \leavevmode%
  \hbox{%
  \begin{beamercolorbox}[wd=.7\paperwidth,ht=2.25ex,dp=1ex,left]{author in head/foot}%
    \hspace*{2ex}\insertshortauthor~~(\insertshortinstitute) \hfill \insertshorttitle
  \end{beamercolorbox}%
  \begin{beamercolorbox}[wd=.3\paperwidth,ht=2.25ex,dp=1ex,right]{date in head/foot}%
    \insertshortdate{}\hspace*{2em}
    \insertframenumber{} / \inserttotalframenumber\hspace*{2ex}
  \end{beamercolorbox}}%
  \vskip0pt%
}

% Configuración de Listings de Python para visualización pedagógica de código
\lstset{
    language=Python,
    basicstyle=\ttfamily\footnotesize,
    keywordstyle=\color{TecAzulOscuro}\bfseries,
    stringstyle=\color{TecRojo},
    commentstyle=\color{TecGris}\itshape,
    showstringspaces=false,
    breaklines=true,
    frame=lines,
    rulecolor=\color{TecAzulOscuro!30},
    backgroundcolor=\color{TecFondoBloque},
    aboveskip=0.5em,
    belowskip=0.5em,
    literate={á}{{\'a}}1 {é}{{\'e}}1 {í}{{\'i}}1 {ó}{{\'o}}1 {ú}{{\'u}}1
             {Á}{{\'A}}1 {É}{{\'E}}1 {Í}{{\'I}}1 {Ó}{{\'O}}1 {Ú}{{\'U}}1
             {ñ}{{\~n}}1 {Ñ}{{\~N}}1 {¿}{{?`}}1 {¡}{{!`}}1
}

% Metadatos institucionales por defecto
\title[Modelación Estadística]{Modelación Estadística para Ciencia de Datos e Ingeniería}
\author[J. Castillo Colmenares]{Juliho David Castillo Colmenares}
\institute[Tec de Monterrey]{Tecnológico de Monterrey \\ Escuela de Ingeniería y Ciencias}
\date{\today}

% Comandos y macros estadísticas compartidas para las presentaciones Beamer (Metropolis)

% Conjuntos numéricos básicos
\newcommand{\R}{\mathbb{R}}
\newcommand{\N}{\mathbb{N}}
\newcommand{\Q}{\mathbb{Q}}
\newcommand{\Z}{\mathbb{Z}}
\newcommand{\set}[1]{\left\{ #1 \right\}}
\newcommand{\abs}[1]{\left|#1\right|}
\newcommand{\norm}[1]{\left\| #1 \right\|}

% Comandos estadísticos y probabilísticos del curso
\newcommand{\Var}{\operatorname{Var}}
\newcommand{\cov}{\operatorname{Cov}}
\newcommand{\corr}{\rho}
\newcommand{\comb}[2]{\begin{pmatrix} #1 \\ #2 \end{pmatrix}}
\newcommand{\s}{\sigma}
\newcommand{\particion}{\mathfrak{P}}
\newcommand{\card}[1]{\#\left( #1 \right)}
\newcommand{\seq}[3]{\set{#1}_{#2}^{#3}}
\newcommand{\E}{\mathbb{E}}
\newcommand{\Prob}{\mathbb{P}}

% Entornos pedagógicos y cajas en español académico natural
\newenvironment{discusion}{%
  \begin{alertblock}{Pregunta de Discusión}%
}{%
  \end{alertblock}%
}

\newenvironment{reflexion}{%
  \begin{alertblock}{Reflexión para el Aula}%
}{%
  \end{alertblock}%
}

\newenvironment{paradoja}{%
  \begin{alertblock}{Paradoja de Exploración}%
}{%
  \end{alertblock}%
}

\newenvironment{motivacion}{%
  \begin{exampleblock}{Motivación: Ciencia de Datos en la Práctica}%
}{%
  \end{exampleblock}%
}

\newenvironment{retoclase}{%
  \begin{block}{Reto Rápido en Equipo}%
}{%
  \end{block}%
}


\title[Contingencia e Independencia]{Sección 07.04: Tablas de Contingencia y Pruebas de Independencia}
\subtitle{Independencia, Homogeneidad, y la Identidad $Z^2=\chi^2$}
\author[J. Castillo Colmenares]{Juliho Castillo Colmenares}
\institute{Tecnológico de Monterrey}
\date{\vspace{-1.2cm}}

\begin{document}

% Slide 1: Portada
\begin{frame}[plain]
  \vspace{-0.8cm}
  \titlepage
\end{frame}

% Slide 2: Hoja de Ruta
\begin{frame}{Hoja de Ruta --- Capítulo 07: Docimasia (Pruebas de Hipótesis)}
  \footnotesize
  ¿Dónde nos encontramos en el desarrollo modular del capítulo? \pause
  \begin{itemize}\itemsep=0.08em
    \item \textbf{07.01} Fundamentos: $H_0$ vs. $H_1$, Errores Tipo I y II, y Potencia \pause
    \item \textbf{07.02} Pruebas $Z$ y $t$ para Medias de Una y Dos Muestras \pause
    \item \textbf{07.03} Pruebas de Bondad de Ajuste $\chi^2$ \pause
    \item \textbf{07.04} Tablas de Contingencia y Pruebas de Independencia \emph{(Hoy --- Cierre del Capítulo)}
  \end{itemize}
  \vspace{-0.05cm}
  \begin{block}{Objetivo de la Sesión}
    Extender el estadístico $\chi^2$ de la Sección 07.03 (una variable categórica) a \textbf{dos} variables categóricas simultáneas: distinguir formalmente entre pruebas de independencia y de homogeneidad, y demostrar la conexión algebraica exacta entre la prueba $\chi^2$ y la prueba $Z$ de dos proporciones en tablas $2\times2$.
  \end{block}
\end{frame}

% Slide 3: Motivación
\begin{frame}{De Una Variable a Dos Variables Categóricas}
  \footnotesize
  \begin{motivacion}
    En la Sección 07.03 preguntamos si \emph{una} variable categórica sigue una distribución teórica. Ahora preguntamos algo distinto y más rico: ¿están \emph{relacionadas} dos variables categóricas? ¿El género influye en la elección de carrera? ¿La generación influye en el modelo de monetización preferido?
  \end{motivacion}

  \vspace{0.15cm}
  \pause
  \begin{itemize}\itemsep=0.1cm
    \item \textbf{Misma fórmula matemática}, $\sum(O-E)^2/E$, en ambos casos. \pause
    \item \textbf{Diferencia crucial:} el diseño experimental --- cómo se recolectó la muestra --- determina si se llama ``independencia'' u ``homogeneidad''.
  \end{itemize}
\end{frame}

% Slide 4: Independencia vs Homogeneidad
\begin{frame}{Independencia vs. Homogeneidad: la Diferencia Conceptual}
  \footnotesize
  \begin{block}{Prueba de Independencia}
    Se extrae \textbf{una sola muestra} de tamaño $N$ de una única población, y se clasifica \emph{a posteriori} según dos variables $A$ y $B$. Ningún total marginal está fijado de antemano (salvo $N$). $H_0: P(A_i\cap B_j)=P(A_i)P(B_j)$.
  \end{block}
  \pause

  \vspace{0.1cm}
  \begin{alertblock}{Prueba de Homogeneidad}
    Se extraen $c$ muestras \textbf{independientes} de tamaños prefijados ($n_1,\ldots,n_c$, columnas fijas por diseño) de $c$ subpoblaciones distintas, midiendo una sola variable categórica. $H_0: p_{i1}=p_{i2}=\ldots=p_{ic}$ para cada categoría $i$.
  \end{alertblock}
\end{frame}

% Slide 5: Estadístico común
\begin{frame}{El Estadístico Común para Ambas Pruebas}
  \footnotesize
  \begin{block}{Frecuencias Esperadas y Estadístico $\chi^2$}
    $$E_{ij} = \frac{R_i C_j}{N}, \qquad \chi^2 = \sum_{i=1}^r\sum_{j=1}^c \frac{(O_{ij}-E_{ij})^2}{E_{ij}} \ \sim\ \chi^2_{(r-1)(c-1)}$$
  \end{block}
  \pause

  \vspace{0.1cm}
  \begin{alertblock}{Punto Clave}
    A pesar de la diferencia conceptual de diseño muestral, la fórmula de cálculo, los grados de libertad $(r-1)(c-1)$, y el procedimiento de decisión son \textbf{idénticos} para independencia y homogeneidad. La distinción importa para la \emph{interpretación}, no para el cálculo.
  \end{alertblock}
\end{frame}

% Slide 6: Identidad Z^2 = chi^2
\begin{frame}{La Identidad Exacta $Z^2=\chi^2$ en Tablas $2\times2$}
  \footnotesize
  \begin{block}{Deducción Algebraica}
    Para una tabla $2\times2$ con celdas $a,b,c,d$: $\chi^2 = \dfrac{N(ad-bc)^2}{(a+b)(c+d)(a+c)(b+d)}$.
  \end{block}
  \pause

  \vspace{0.1cm}
  \begin{alertblock}{Conexión con la Prueba $Z$ de Dos Proporciones}
    Como $Z\sim N(0,1) \implies Z^2\sim\chi^2_1$, y la prueba $Z$ de dos proporciones (Sección 07.02) es estructuralmente la raíz cuadrada firmada del estadístico de homogeneidad en una tabla $2\times2$: se cumple exactamente $Z^2\equiv\chi^2$. Ambas docimasias son matemáticamente idénticas para el caso binario.
  \end{alertblock}
\end{frame}

% Slide 7: Aplicaciones
\begin{frame}{Aplicaciones en Ciencia de Datos e Ingeniería}
  \footnotesize
  \begin{itemize}\itemsep=0.1cm
    \item \textbf{Segmentación de mercado:} ¿el modelo de monetización preferido depende de la generación del usuario? \pause
    \item \textbf{Ensayos clínicos:} ¿la tasa de infección depende del brazo de tratamiento (vacuna vs. placebo)? \pause
    \item \textbf{Auditoría de equidad algorítmica:} ¿las decisiones de un modelo son independientes de un atributo protegido? \pause
    \item \textbf{Comparación multi-grupo (Marascuilo):} identificar exactamente qué pares de grupos difieren tras un rechazo global.
  \end{itemize}
\end{frame}

% Slide 8: Lab Python (Bloque 1)
\begin{frame}[fragile]{Laboratorio en Python: Prueba de Independencia}
  \scriptsize
  Tabla $4\times3$ (tipo de vehículo vs. día), verificada contra \texttt{scipy.stats.chi2\_contingency} (\texttt{07.04\_contingency\_tables.py}):
  {\lstset{basicstyle=\fontsize{3.8pt}{4.5pt}\selectfont\ttfamily,aboveskip=0.05em,belowskip=0.05em}
  \lstinputlisting[firstline=17, lastline=32]{../../code/07_pruebas_hipotesis/07.04_contingency_tables.py}}
\end{frame}

% Slide 9: Lab Python (Bloque 2)
\begin{frame}[fragile]{Laboratorio en Python: Prueba de Homogeneidad}
  \scriptsize
  Comparación de tres cohortes generacionales independientes, misma fórmula distinto diseño:
  {\lstset{basicstyle=\fontsize{4pt}{4.8pt}\selectfont\ttfamily,aboveskip=0.1em,belowskip=0.1em}
  \lstinputlisting[firstline=43, lastline=64]{../../code/07_pruebas_hipotesis/07.04_contingency_tables.py}}
\end{frame}

% Slide 10: Lab Python (Bloque 3)
\begin{frame}[fragile]{Laboratorio en Python: Identidad $Z^2=\chi^2$}
  \scriptsize
  Verificación numérica exacta de la equivalencia algebraica en una tabla $2\times2$:
  {\lstset{basicstyle=\fontsize{3.8pt}{4.5pt}\selectfont\ttfamily,aboveskip=0.05em,belowskip=0.05em}
  \lstinputlisting[firstline=67, lastline=82]{../../code/07_pruebas_hipotesis/07.04_contingency_tables.py}}
\end{frame}

% Slide 11: Salida Terminal
\begin{frame}[fragile]{Laboratorio en Python: Salida en Terminal}
  \scriptsize
  Salida estándar generada al ejecutar el script del laboratorio en Python:
  \vspace{0.02cm}
  \begin{block}{Salida Estándar en Consola}
    \fontsize{4pt}{5pt}\selectfont
    \begin{verbatim}
=== Block 1: Independence Test ===
chi2 = 8.4896, df=6, p-value = 0.2044

=== Block 2: Homogeneity Test (3 Cohorts) ===
chi2 = 56.1485, df=4, p-value = 0.000000

=== Block 3: Z^2 = Chi^2 Identity ===
chi2 = 7.3846, Z^2 = 7.3846, |chi2-Z^2| = 0.0000000000
    \end{verbatim}
  \end{block}
\end{frame}

% Slide 16: Cierre
\begin{frame}{Cierre de Sección y de Capítulo: Contingencia e Independencia}
  \begin{reflexion}
    Con esta sección cerramos el Capítulo 07: hemos construido el marco lógico completo de la docimasia --- desde el planteamiento de $H_0/H_a$ y el balance $\alpha$-$\beta$, pasando por medias (Z/t), hasta modelos categóricos completos ($\chi^2$ de bondad de ajuste y de contingencia). Toda inferencia estadística aplicada descansa sobre estas herramientas.
  \end{reflexion}

  \vspace{0.2cm}
  \pause
  \begin{block}{Lecciones Clave}
    \begin{itemize}\itemsep=0.08cm
      \item \textbf{Independencia vs. homogeneidad:} misma fórmula, diseño muestral distinto.
      \item \textbf{Estadístico común:} $\chi^2=\sum(O_{ij}-E_{ij})^2/E_{ij}$, $\nu=(r-1)(c-1)$.
      \item \textbf{Identidad $Z^2=\chi^2$:} las pruebas de proporciones y de contingencia $2\times2$ son matemáticamente idénticas.
    \end{itemize}
  \end{block}
\end{frame}

% Slide 17: Perspectiva Modular
\begin{frame}{Perspectiva Modular --- El Próximo Reto}
  \scriptsize
  \begin{retoclase}
    Con el Capítulo 07 completo, dominamos la comparación de \textbf{dos} grupos o variables. El Capítulo 08 (Diseño de Experimentos) generaliza esta lógica a la comparación simultánea de \textbf{tres o más} grupos mediante el Análisis de Varianza (ANOVA) de un factor y la prueba $F$, evitando la inflación del error Tipo I que surgiría de múltiples pruebas $t$ por pares.
  \end{retoclase}

  \vspace{0.08cm}
  \pause
  \begin{alertblock}{Preparación Recomendada}
    \begin{itemize}\itemsep=0.06cm
      \item Repasa la prueba $F$ de Fisher-Snedecor y su interpretación como cociente de varianzas (Sección 05.05).
      \item Reflexiona sobre por qué realizar $\binom{k}{2}$ pruebas $t$ por pares infla el error Tipo I familiar (FWER).
    \end{itemize}
  \end{alertblock}
\end{frame}

\end{document}
